\begin{table}[htbp]
\centering
\begin{threeparttable}
\caption{Concentration of relationships by segment, 2017--2023}
\label{tab:concentration}
\begin{tabular}{lrrr}
\toprule
 & Hardware & Software & IT services \\
 & (CPV 30) & (CPV 48) & (CPV 72) \\
\midrule
\multicolumn{4}{l}{\emph{Share of all relationships held by the largest suppliers (\%)}} \\
\quad Largest supplier & 6.1 & 3.7 & 1.6 \\
\quad Five largest & 15.7 & 13.5 & 6.1 \\
\quad Ten largest & 22.7 & 20.3 & 9.5 \\
\addlinespace
\multicolumn{4}{l}{\emph{Share of all relationships held by the largest buyers (\%)}} \\
\quad Largest buyer & 0.9 & 1.6 & 2.5 \\
\quad Five largest & 3.2 & 6.2 & 8.9 \\
\quad Ten largest & 5.2 & 10.4 & 13.7 \\
\addlinespace
\multicolumn{4}{l}{\emph{Segment size}} \\
\quad Distinct relationships & 7,917 & 5,888 & 9,157 \\
\quad Suppliers & 2,250 & 2,355 & 4,257 \\
\quad Buyers & 2,353 & 2,013 & 2,116 \\
\bottomrule
\end{tabular}
\begin{tablenotes}[flushleft]
\footnotesize
\item \textit{Note.} A relationship is a distinct buyer--supplier pair observed at least once between 2017 and 2023; repeated awards within a pair are counted once. Shares on both sides are expressed as a proportion of the same total, so supplier and buyer concentration are directly comparable. All buyers and suppliers with at least one award in the segment are included; no degree filter is applied.
\end{tablenotes}
\end{threeparttable}
\end{table}
